%% Section 5: N\'{e}ron--Tate height and height on the base


We keep the notation of $\mathsection$\ref{SectionSettingUpHtIneq}. In
particular, $S$ is a regular, irreducible, quasi-projective variety over $\IQbar$ and
$\pi\colon\cA\rightarrow S$ is an abelian scheme of relative dimension
$g \ge 1$. Moreover, we have immersions as in
$\mathsection$\ref{SectionSettingUpHtIneq} and we
assume {\tt (Hyp)} from page~\pageref{def:hyp}.
We use the heights introduced in $\mathsection$\ref{SubsectionHeight}. 
Let $X$ be an irreducible closed subvariety of $\cA$ 
defined over $\IQbar$. We assume that $X$ dominates $S$ and is
non-degenerate as defined in Definition~\ref{DefinitionNonDegenerate}

The upshot of {\tt (Hyp)} is that we obtain from Proposition~\ref{PropBettiForm} the Betti form $\omega$ on $\an{\cA}$. Moreover, part (i) and (iii) of Proposition~\ref{PropBettiForm} implies that, for $d = \dim X$, 
\begin{equation}\label{EqBettiFormHtIneq}
\omega|_{X^{\mathrm{sm},\mathrm{an}}}^{\wedge d} > 0 \quad \text{ at some smooth point of }X^{\mathrm{an}}.
\end{equation}


Our assumption \eqref{EqBettiFormHtIneq} allows us to apply
Proposition~\ref{PropAuxHtIneq} to
$X$. 
 There exists a constant $c_1>0$ as in \eqref{EqConstantInHtIneq} such
 that the following holds. Let $N\in\IN$ be a power of $2$, there
 exists a Zariski open dense subset $U_N\subset X$ and a constant
 $c_2(N)\ge 0$ such that
\begin{equation}\label{EquationHeight1}
h([N]P) \ge c_1 N^2 h(\pi(P)) - c_2(N)
\end{equation}
for all $P\in U_N(\IQbar)$; we stress that $U_N$ and 
 $c_2\ge 0$ may depend  on $N$ in addition to $X,\cA,$ and the various
 immersions such as $\cA \subseteq \IP_\IQbar^n \times \IP_\IQbar^m$.


By the  Theorem of Silverman-Tate, see \cite[Theorem~A]{Silverman} and
Theorem~\ref{thm:silvermantate}, there exist a  constant
$c_0\ge 0$ %depending only on $\mathcal{A}/S$ and 
 such that
\begin{equation}
\label{EquationHeight2}
|\hat{h}_{\mathcal{A}}(P)-h(P)| \le
c_0\max\{1,h(\pi(P))\}\le c_0\bigl(1+h(\pi(P))\bigr)
% c_3 h_{\overline S,\mathcal{M}}(\pi(P))+c_4
\end{equation}
for all $P\in \mathcal{A}(\IQbar)$.


Next we kill Zimmer constants as in
Masser's~\cite[Appendix~C]{ZannierBook}.
For any $P\in U_N(\IQbar)$, we have
\begin{align*}
\hat{h}_{\cA}([N](P)) &\ge h([N](P))
-c_0\bigl(1+h(\pi([N](P)))\bigr) 
&&\text{(by (\ref{EquationHeight2}))}\\
&= h([N](P))
-c_0\bigl(1+h(\pi(P))\bigr)
&&\text{(as $\pi([N](P))=\pi(P)$)}
\\
&\ge c_1 N^2 h(\pi(P)) -c_2(N) -c_0\bigl(1+h(\pi(P))\bigr)
&&\text{(by (\ref{EquationHeight1})).}
\end{align*}
We use $\hat{h}_{\mathcal{A}}([N]P)=N^2 \hat{h}_{\mathcal{A}}(P)$, 
divide by $N^2$, and rearrange to get
\begin{equation*}%\label{EquationFinal}
\hat{h}_{\cA}(P) \ge \left(c_1 - \frac{c_0}{N^2} \right) h(\pi(P)) - \frac{c_2(N)+ c_0}{N^2}
\end{equation*}
for all $N\in\IN$ that are powers of $2$  and all $P \in U_N(\IQbar)$. 

Recall that $c_0$ and $c_1$ are independent of $N$. We  fix $N\in\IN$ to be the least
power of $2$ such that $N^2\ge 2c_0/c_1$. As $h(\pi(P))$ is
non-negative we get  % that \eqref{EquationFinal} holds for $N'$, then
\[
\hat{h}_{\cA}(P) \ge \frac{c_1}{2} h(\pi(P)) - \frac{c_2(N)+ c_0}{N^2}
\]
for all
$P \in U_N(\IQbar)$. Since $N$ is now fixed, the Zariski open dense
subset $U_N$ of $X$ is also fixed.
The theorem follows after adjusting $c_1$ and $c_2$. 
\qed
%\end{proof}
\begin{remark}
In the proof of Theorem~\ref{ThmHtInequality} we can keep track of the
process to compute the constant $c_1 > 0$. Use the notation in
$\mathsection$\ref{SectionSettingUpHtIneq}. In particular $\omega$ is
the Betti form on $\cA$, we have an immersion $\cA \subseteq \IP_\IC^n
\times \IP_\IC^m$, $\alpha$ is a $(1,1)$-form on $\IP_\IC^n \times
\IP_\IC^m$ representing the Chern class of $\cO(1,1)$, $\Delta \subseteq S^{\mathrm{an}}$ is open and relative compact, and $\theta \colon \cA^{\mathrm{an}} \rightarrow [0,1]$ (which factors through $S^{\mathrm{an}}$) is a smooth function with compact support contained in $\cA_{\Delta} := \pi^{-1}(\Delta)$. The function $\theta$ should furthermore satisfy $\theta(P_0) = 1$ for some $P_0 \in X^{\mathrm{sm}}(\IC)$ such that $(\omega|_{X^{\mathrm{sm}}})^{\wedge d}$ is positive at $P_0$, where $d = \dim X$.

Assume $d \ge 1$. The proof of Theorem~\ref{ThmHtInequality} tells us
that one half of any rational number satisfying the inequality
\eqref{EqConstantInHtIneq} can be taken as $c_1$. So the constant $c_1 > 0$ can be taken to be any rational
number
in  $(0,\kappa/(2c d))$, such that: %where  $\kappa$ and $c$ are:
\begin{itemize}
\item $\kappa = \kappa'/C^d$, where $\kappa' = \int_{\sman{X}}(
  \theta\omega )^{\wedge d}$, as in (\ref{eq:N2dgrowth}), and $C$ satisfies $C
  \alpha|_{\cA_{\Delta}} - \omega|_{\cA_{\Delta}} \ge 0$, as in (\ref{EqConditionOnCapitalC}),
\item $c$ is a  constant depending on a certain degree of  $X$ and coming from
  (\ref{eq:intersectionnumberub2}). % $c' = D^{\prime d-1}\binom{n+d}{n} (\cO(1,1)^{\cdot d}[\overline{X}])$, where $D'$ is such that the bihomogeneous polynomials defining $[2] \colon \cA \rightarrow \cA$ have bi-degree $(D',4)$.
\end{itemize}
\end{remark}

%%% Local Variables:
%%% TeX-master: "main"
%%% End:
